\documentclass[10pt,a4paper]{amsart}
\usepackage[margin=19mm]{geometry}
\usepackage{amsmath,amssymb,mathtools}
\usepackage[scr=rsfs,cal=euler]{mathalfa}
\usepackage{xfrac}
\usepackage[unicode,hidelinks,bookmarks=false]{hyperref}
\newcommand{\ie}{i.e.\ }
\newcommand{\cf}{cf.\ }
\newcommand{\resp}{respectively\ }
\newcommand{\Subsection}{\subsection}
\providecommand{\nobreakdash}{\nobreak\hbox{-}}
\setlength{\emergencystretch}{2em}
\multlinegap0pt
\usepackage{bm}
\usepackage{bbm}
\usepackage{dsfont}
\usepackage{xspace}
\usepackage{cancel}
\usepackage{mathtools}
\usepackage{amsthm}
\makeatletter
\@ifundefined{lemma}{\newtheorem{lemma}{Lemma}}{}
\@ifundefined{theorem}{\newtheorem{theorem}{Theorem}}{}
\makeatother
\title[Quantum Advantage in Learning Mixed Unitary Channels]{Quantum Advantage in Learning Mixed Unitary Channels}
\author{Yue Tu, Liang Jiang}
\date{Source: arXiv:2511.13683v1 (CC BY 4.0)}
\begin{document}
\maketitle

\noindent\emph{Source and licence.} Quantum Advantage in Learning Mixed Unitary Channels. Version arXiv:2511.13683v1 (CC BY 4.0), distributed under the Creative Commons Attribution 4.0 International licence, as recorded in the arXiv licence metadata for this exact version and shown on the abstract page https://arxiv.org/abs/2511.13683v1. Full rights notice: licenses/arxiv-2511.13683-CC-BY-4.0.txt.\par\smallskip

\noindent\textit{Editorial scope.} Counted unit: the Theorem whose source label is \texttt{th:trFI} in the article (source0.tex lines 349--358, source label \texttt{th:trFI}), delivered together with the article's own complete original proof of it (source0.tex lines 360--411, one \texttt{proof} environment of 52 lines). No mathematics is written, repaired or re-derived: statement and proof are the article's own text line for line. Required context retained uncounted: eq:FIM (lines 197--199); eq:concFI (lines 277--280); lm:suc (lines 329--334). Every definition, display and label of those lines is retained unchanged. Omitted from this package and not redistributed: 1--196, 200--276, 281--328, 335--348, 359--359, 412--700. Nothing inside a retained excerpt is omitted or altered: every retained body is delivered exactly as the article's lines are written, so no omission receipt of any kind accompanies this package. Citations: the retained excerpts carry no citation command at all, so no bibliography entry of the article is transcribed and no reference is redistributed. Reference adaptation: none was needed and none was made; every reference inside the delivered unit resolves inside it, so no unresolved reference or citation is produced. Printed slips kept literal: no printed word, symbol, hypothesis or number of the counted statement or of its proof was altered. Layout and class: the article's own class is \texttt{article} with packages including geometry, amsmath, amsthm, amssymb, bm, braket, graphicx, algorithm, algpseudocode, hyperref; the wrapper is the lane's pinned \texttt{amsart} editorial wrapper at 10pt on A4, which restates the theorem-like environments the retained text needs and adds the article's own macro definitions that the retained text needs (none), with extra wrapper packages (bm, bbm, dsfont, xspace, cancel, mathtools). The delivered numbering is the unit's own. Assets: no retained span contains a figure environment, a graphics-inclusion command, a PDF-page inclusion, a table or a TikZ picture; no image file is copied into this package, the package declares no asset dependency and no image file is redistributed with it. Licence: version arXiv:2511.13683v1 of this article is distributed under the Creative Commons Attribution 4.0 International licence, as recorded in the arXiv licence metadata for this exact version and shown on the abstract page https://arxiv.org/abs/2511.13683v1. A full rights notice is delivered at licenses/arxiv-2511.13683-CC-BY-4.0.txt. Characters: every retained excerpt is pure ASCII, so the delivered engine, XeLaTeX with its default Latin Modern text font, reports no missing character.
\begin{equation}\label{eq:FIM}
    I(\theta) = P_s K^\top D(p)^{-1} K P_s,
\end{equation}

\begin{equation}\label{eq:concFI}
    I^{(k)}(\theta)
    = J(\theta)^{\top} \, \tilde{I}(\tilde{\theta}) \, J(\theta),
\end{equation}

\begin{lemma}\label{lm:suc}
For the $K$ matrix in Eq.~\ref{eq:FIM}, we have
\begin{equation}
    \sum_{i=1}^s \max_j \{K_{ij}\} \leq d.
\end{equation}
\end{lemma}

\begin{theorem}\label{th:trFI}
For the non-concatenating protocol, we have
\begin{equation}\label{eq:trInonc}
    \operatorname{Tr}(I(u)) \le r d.
\end{equation}
For the concatenating protocol with at most $k$ concatenations, we have
\begin{equation}\label{eq:trIc}
    \operatorname{Tr}(I(u)) \le k^2 r d.
\end{equation}
\end{theorem}

\begin{proof}
We first prove Eq.~\ref{eq:trInonc}. From Eq.~\ref{eq:FIM}, we have
\begin{equation}
\begin{split}
    \operatorname{Tr}I(u)
    &= \operatorname{Tr}(P_s K^\top D(p)^{-1} K P_s) \\
    &\le \operatorname{Tr}(K^\top D(p)^{-1} K) \\
    &= \sum_{t=1}^{s}\frac{\sum_{j=1}^{r} K_{tj}^{2}}{p_t} \\
    &\le \sum_{t=1}^{s}\frac{(\max_{1\le j\le r} K_{tj})\sum_{j=1}^{r} K_{tj}}{p_t} \\
    &= \sum_{t=1}^{s}\frac{(\max_{1\le j\le r} K_{tj})\sum_{j=1}^{r} K_{tj}}{[Ku]_t} \\
    &= \sum_{t=1}^{s}\frac{(\max_{1\le j\le r} K_{tj})\sum_{j=1}^{r} K_{tj}}{\frac{1}{r}\sum_{j=1}^{r} K_{tj}} \\
    &= r\sum_{t=1}^{s}\max_{1\le j\le r} K_{tj}\\
    &\le r\,d.
\end{split}
\end{equation}
Where the final inequality uses Lemma~\ref{lm:suc}. Thus, the bound in Eq.~\ref{eq:trInonc} is established.

Now we prove Eq.~\ref{eq:trIc}. From Eq.~\ref{eq:concFI}, we have
\begin{equation}\label{eq:trIJconc}
    \operatorname{Tr}(I^{(k)}(\theta)) 
    = \operatorname{Tr}(J(\theta)^{\top} \, \tilde{I}(\tilde{\theta}) \, J(\theta)) 
    \le \operatorname{Tr}(\tilde{I}(\tilde{\theta})) \| J(\theta) \|_2^2 
    \le r^k d  \| J(\theta) \|_2^2,
\end{equation}
where the second inequality follows because $\tilde{I}(\tilde{\theta})$ is the expectation of rank-one operators, and the last inequality because the effective mixed unitary channel has rank $r^k$.

Next, we upper bound $\| J(\theta) \|_2^2$.
For the tensor map $f(\theta) = \theta^{\otimes k}$ with $\theta \in \mathbb{R}^r$, 
the Jacobian at the uniform point $\theta_* = \mathbf{1}/r$ has entries
\begin{equation}
    J_{ab} = m(a,b)\,r^{1-k},
\end{equation}
where $m(a,b)$ counts how many times index $b$ appears in the $k$-tuple $a \in [r]^k$.
A direct combinatorial calculation gives
\begin{equation}
    J^\top J = k\,r^{-k}\!\left(I_r + (k-1)\mathbf{1}\mathbf{1}^\top\right),
\end{equation}
whose largest eigenvalue is
\begin{equation}
    \lambda_{\max} = k\,r^{-k}\!\left(1 + (k-1)r\right).
\end{equation}
Hence, the operator two-norm of $J$ is
\begin{equation}
    \|J(\theta_*)\|_2
    = \sqrt{k\,r^{-k}\!\left(1 + (k-1)r\right)} \le k r^{(1-k)/2}.
\end{equation}
Substituting this into Eq.~\ref{eq:trIJconc}, we obtain
\begin{equation}
    \operatorname{Tr}(I^{(k)}(\theta_*)) \le r^k d  \| J(\theta) \|_2^2 \le k^2 r d,
\end{equation}
which completes the proof of Eq.~\ref{eq:trIc}.
\end{proof}

\end{document}
